\subsection{有限生成扩张的微分性质}

定理 4. --- 设 P 为完全域，L 为 P 的扩域，K 为 L 的子扩张；我们假设 L 是 K 的有限生成扩张。则典范 L-线性映射 \(\alpha: \Omega_{\mathrm{P}}(\mathrm{K}) \otimes_{\mathrm{K}} \mathrm{L} \to \Omega_{\mathrm{P}}(\mathrm{L})\) 的指标等于 L 在 K 上的超越次数。

如果 E 和 F 是 L 的两个子扩张，使得 \(E \subset F\) ，我们用 \(\alpha(F/E)\) 表示 \(\Omega_{\mathrm{P}}(\mathrm{E}) \otimes_{\mathrm{E}} F\) 到 \(\Omega_{\mathrm{P}}(\mathrm{F})\) 中的典范 F-线性映射，当它有定义时， \(\alpha(F/E)\) 的指标将记作 \(d(F/E)\) 。如果 E、F 和 G 是 L 的三个子扩张，使得 \(E \subset F \subset G\) ，我们有一个交换图

\[
\begin{array}{c} \Omega_ {\mathrm{P}} (\mathrm{F}) \otimes_ {\mathrm{F}} \mathrm{G} \xrightarrow {\alpha (\mathrm{G} / \mathrm{F})} \Omega_ {\mathrm{P}} (\mathrm{G}) \\ \alpha (\mathrm{F} / \mathrm{E}) \otimes_ {\mathrm{F}} \mathrm{Id} _ {\mathrm{G}} \Bigg | _ {\uparrow} \\ (\Omega_ {\mathrm{P}} (\mathrm{E}) \otimes_ {\mathrm{E}} \mathrm{F}) \otimes_ {\mathrm{F}} \mathrm{G} \xrightarrow {\beta} \Omega_ {\mathrm{P}} (\mathrm{E}) \otimes_ {\mathrm{E}} \mathrm{G} \end{array}
\]

其中 \(\beta\) 是命题 2（II，第 278 页）中定义的典范同构。由于指标在标量扩张下显然不变，且同构的指标为零，引理 2（V，第 133 页）表明指标 \(d(\mathrm{G}/\mathrm{E})\) 在 \(d(\mathrm{F}/\mathrm{E})\) 与 \(d(\mathrm{G}/\mathrm{F})\) 皆有定义时有定义，并且此时

\[
d (\mathrm{G} / \mathrm{E}) = d (\mathrm{G} / \mathrm{F}) + d (\mathrm{F} / \mathrm{E}).\tag{5}
\]

由于 L 是 K 的有限生成扩张，公式 (5) 和 V，第 111 页的推论表明，只需考虑存在 x 使得 \(L = K(x)\) 的情形；进一步地，若 x 在 K 上是代数的，则存在 L 的特征指数的某个幂 q，使得 \(x^{q}\) 在 K 上是可分代数的（V，第 44 页，命题 13）。因此，只需在以下三种特殊情形下建立等式 \(d(\mathrm{L}/\mathrm{K}) = \operatorname{tr} \cdot \deg_{\mathrm{K}} \mathrm{L}\) ：

\begin{enumerate}
\def\labelenumi{\alph{enumi})}
\item
  \(x\) 在 \(\mathbf{K}\) 上是超越的：则 \(\alpha\) 是单射且其余核在 \(\mathrm{L}(\mathrm{V},\mathrm{p}.129,\mathrm{Cor.})\) 上的秩为 1；所以我们有 \(d(\mathrm{L} / \mathrm{K}) = 1\) 并且 tr. \(\deg_{\mathbf{K}}\mathrm{L} = 1\) 。
\item
  \(x\) 在 \(\mathbf{K}\) 上是可分代数的：则 \(\alpha\) 是双射（V，第 130 页，推论 2），从而 \(d(L / K) = 0\) ；显然也有 tr. \(\deg_{\mathbf{K}}\mathbf{L} = 0\) 。
\item
  域 L 的特征为 \(p \neq 0\) , \(x \notin K\) 且 \(x^{p} = a\) 属于 K： \(\alpha\) 的余核 C 同构于 \(\Omega_{\mathrm{K}}(\mathrm{L})\) ，且由于 \(\{x\}\) 是 L 在 K 上的 p-基，空间 C 在 L 上的维数为 1（V，第 102 页，定理 1）。由于 a 是 L 中的 p 次幂，我们有 \(d_{L/P}a = 0\) 且 \(\alpha\) 的核包含 \(U = \Omega_{\mathrm{P}}(\mathrm{K}) \otimes_{\mathrm{K}} L\) 的由 \(d_{K/P}a \otimes 1\) 生成的子空间 R。对于每个 \(y \in K\) ，令 \(\Delta(y)\) 为 \(d_{K/P}y \otimes 1 \mod R\) 的剩余类；则 \(\Delta\) 是 K 到 U/R 中的一个 P-导子，使得 \(\Delta(a) = 0\) 。命题 5（V，第 101 页）表明 \(\Delta\) 可延拓为 L 到 U/R 中的一个 P-导子 D。因此存在一个 L-线性映射 \(\beta : \Omega_{\mathrm{P}}(\mathrm{L}) \to \mathrm{U}/\mathrm{R}\) 使得 \(D = \beta \circ d_{L/P}\) 且 \(\beta \circ \alpha\) 是 U 到 U/R 上的典范映射。这就证明了 R 是 \(\alpha\) 的核。由于 P 是完全域，我们有 \(P(K^{p}) = K^{p}\) ，从而 \(a \notin P(K^{p})\) ，最终 \(d_{K/P}a \neq 0\) （V，第 103 页，命题 6）。于是 \(\alpha\) 的核与余核的维数均为 1，从而 \(d(\mathrm{L}/\mathrm{K}) = 0\) ，并且我们也有 tr. deg \(_{K}L = 0\) 。
\end{enumerate}

推论 1. --- 设 L 是域 K 的有限生成扩张，其超越次数为 s。向量空间 \(\Omega_{\mathrm{K}}(\mathrm{L})\) 在 L 上的维数为 \(\geqslant s\) ，当且仅当 L 在 K 上可分时取等号。

设 P 为 K 的素子域。设 N 为 \(\alpha\) 的核，则由序列 \((\mathrm{E}_{\mathrm{P},\mathrm{K},\mathrm{L}})(\mathrm{V},\mathrm{p}.128)\) 的正合性及定理 4，我们有 \([\Omega_{\mathrm{K}}(\mathrm{L}):\mathrm{L}] = s + [\mathrm{N}:\mathrm{L}]\) ；由 V，第 132 页，推论，K 的扩张 L 是可分的当且仅当 N = 0。推论随之成立。

推论 2. --- 设 L 是域 K 的有限生成扩张。要使 L 在 K 上是代数的且可分的，其充分必要条件为 \(\Omega_{\mathbf{K}}(\mathbf{L}) = 0\) 。

这直接由推论 1 得出。

推论 3. --- 设 K 是一个特征为 \(p \neq 0\) 的域，L 是 K 的超越次数为 s 的有限生成扩张。若 \([K:K^{p}]\) 是有限的，则 \([L:L^{p}]\) 亦然，并且我们有 \([L:L^{p}] = p^{s} \cdot [K:K^{p}]\) 。

设 P 为 K 的素子域。若 \([K:K^{p}]\) 是有限的，则向量空间 \(\Omega_{\mathrm{P}}(\mathrm{K})=\Omega_{\mathrm{K}^{\mathrm{p}}}\left(\mathrm{K}\right)\) 在 K 上具有有限维数 m，并且我们有 \([K:K^{p}]=p^{m}(V,p.103,Th.2)\) ；于是向量空间 \(\Omega_{\mathrm{P}}(\mathrm{K})\otimes_{\mathrm{K}}L\) 在 L 上的维数也是有限的 m。进一步，由于 K 在 \(K^{p}\) 上的次数有限，域 \(\mathrm{K}(L^{p})\) 在 \(\mathrm{K}^{p}(L^{p})=L^{p}\) 上是有限次的；由于域 L 是 K 的有限生成扩张且在 \(\mathrm{K}(L^{p})\) 上是代数的，它是 \(\mathrm{K}(L^{p})(\mathrm{V},p.18,Th.2)\) 的有限次扩张；于是我们得出结论，L 在 \(L^{p}(V,p.10,Th.1)\) 上是有限次的。于是 \(\Omega_{\mathrm{P}}(\mathrm{L})\) 是 L 上有限维数为 n 的向量空间，并且我们有 \([L:L^{p}]=p^{n}(V,p.103,Th.2)\) 。根据引理 1（V，第 132 页），L-线性映射 \(\alpha:\Omega_{\mathrm{P}}(\mathrm{K})\otimes_{\mathrm{K}}L\to\Omega_{\mathrm{P}}(\mathrm{L})\) 因此具有指标 n-m，由此根据定理 4（V，第 133 页）得 n-m=s，且 \(p^{n}=p^{s}\cdot p^{m}\) 。

注记. --- 1) 设 K 为一个特征为 \(p \neq 0\) 的域，L 是 K 的一个扩域。我们有 \(\Omega_{\mathbf{K}}(\mathbf{L}) = 0\) 当且仅当 \(\mathbf{L} = \mathbf{K}(\mathbf{L}^p)\) （V，第 103 页，命题 6）。因此，若 L 在 K 上有限生成，则 L 是 K 上的代数可分扩张当且仅当我们有 \(\mathbf{L} = \mathbf{K}(\mathbf{L}^p)\) 。当 L 不是在 K 上有限生成时，此结果一般不再成立，正如 L 是 K 的完全闭包的情形所示。

\begin{enumerate}
\def\labelenumi{\arabic{enumi})}
\setcounter{enumi}{1}
\tightlist
\item
  设 K 为一个域， \(F_{1}, \ldots, F_{m}\) 为 \(K[X_{1}, \ldots, X_{n}]\) 中的多项式，且 L 是由元素 \(x_{1}, \ldots, x_{n}\) 生成的 K 的扩域。假设多项式 \(F_{1}, \ldots, F_{m}\) 生成 \(K[X_{1}, \ldots, X_{n}]\) 中由所有满足 \(F(x_{1}, \ldots, x_{n}) = 0\) 的多项式 F 组成的理想。由命题 1（V，第 127 页）和微分模的泛性质（III，第 569 页），我们容易推出以下结果：向量空间 \(\Omega_{K}(L)\) （在 L 上的）由 \(dx_{1}, \ldots, dx_{n}\) 生成；我们有关系
\end{enumerate}

\[
\sum_ {i = 1} ^ {n} \mathrm{D} _ {i} \mathrm{F} _ {j} (x _ {1}, \dots , x _ {n}) \cdot d x _ {i} = 0 \quad (\text { for } 1 \leqslant j \leqslant m);\tag{6}
\]

最后，如果 \(u_1, \ldots, u_n\) 是 \(L\) 的元素，使得 \(\sum_{i=1}^{n} u_i \cdot dx_i = 0\) ，则存在元素 \(v_1, \ldots, v_m\) （\(L\) 的），使得 \(u_i = \sum_{j=1}^{m} D_i F_j(x_1, \ldots, x_n) v_j\) （对于 \(1 \leqslant i \leqslant n\) ）。令 \(r\) 为矩阵 \((D_i F_j(x_1, \ldots, x_n))\) （具有 \(n\) 行和 \(m\) 列）的秩；令 \(s\) 为 L 在 K 上的超越次数。那么我们有 \([\Omega_{\mathrm{K}}(\mathrm{L}):\mathrm{L}]=n-r\) 。因此，K 的扩域 L 是可分的当且仅当 \(r+s=n\) （推论 1），并且它是代数且可分的当且仅当 r=n（推论 2）。

